\section{文献判据：CRLB、GDOP 与几何稀释}\label{sec:lit}

本节把纯方位测向定位的经典理论判据实现成一层可对照的判据，用于\textbf{解释趋势、快速预筛、
独立认证}，并明确其与本文集员判据的边界。本文逐项实现并与数值对账的文献为
Foy~\cite{foy1976}、Chan 与 Ho~\cite{chan1994}、Chen 等~\cite{chen2009} 以及任叶童~\cite{ren2016} 四篇；本项目文献清单中
其余相关条目的原文未随项目归档，故仅在定性趋势层面转述，不逐条引用其公式。

\subsection{两站纯方位测向的 Fisher 信息与 CRLB}

设两站 $S_1,S_2$ 到源的距离为 $R_1,R_2$，观测模型为
$\psi_i=\varphi_i(G)+e_i$，$\varphi_i=\angle(S_i\!\to\!G)$，$e_i$ 独立同分布、标准差 $\sigma$。
记视线单位向量 $L_i=(G-S_i)/R_i$、其法向 $n_i=(-L_{i,y},L_{i,x})^\top$，则单个测向对源的
Fisher 信息为 $\dfrac{1}{\sigma^2}\dfrac{1}{R_i^2}n_in_i^\top$，两站独立叠加得

\begin{equation}
  H=\frac{1}{\sigma^2}\left(\frac{n_1n_1^\top}{R_1^2}+\frac{n_2n_2^\top}{R_2^2}\right),
  \qquad
  \operatorname{tr}H=\frac{1}{\sigma^2}\Bigl(\frac{1}{R_1^2}+\frac{1}{R_2^2}\Bigr),\quad
  \det H=\frac{\sin^2\gamma}{\sigma^4R_1^2R_2^2},
  \label{eq:fim}
\end{equation}

其中 $\gamma$ 为源处两条示向度的交会角。对 $2\times2$ 对称矩阵有
$\operatorname{tr}(H^{-1})=\operatorname{tr}H/\det H$（因 $\operatorname{tr}(\operatorname{adj}H)
=\operatorname{tr}H$），于是位置误差的均方尺度（几何精度因子）为

\begin{equation}
  \mathrm{GDOP}=\sqrt{\operatorname{tr}(H^{-1})}
  =\frac{\sigma\sqrt{R_1^2+R_2^2}}{\sin\gamma}.
  \label{eq:gdop}
\end{equation}

$H$ 的特征值给出 CRLB 误差椭圆的两个半轴（$1\sigma$），其\textbf{主半轴}与\textbf{面积}有闭式：

\begin{equation}
  a=\frac{\sqrt2\,\sigma R_1R_2}
  {\sqrt{\,R_1^2+R_2^2-\sqrt{(R_1^2+R_2^2)^2-4R_1^2R_2^2\sin^2\gamma}\,}},
  \qquad
  S_{\mathrm{ellipse}}=\frac{\pi\sigma^2R_1R_2}{\sin\gamma}.
  \label{eq:crlb}
\end{equation}

式~\eqref{eq:crlb} 与任叶童~\cite{ren2016} 给出的双站误差椭圆半轴式与面积式一致。程序对
$53$ 组几何同时用闭式与数值特征分解（\texttt{numpy.linalg.eigh}）计算，两者最大相对偏差
$1.41\times10^{-12}$，可认为闭式实现正确。面积式最能说明问题：$S_{\mathrm{ellipse}}\propto
R_1R_2/\sin\gamma$，即\textbf{"靠近源"与"张角大"两个目标相互竞争}；Chen 等~\cite{chen2009} 对多观测器
方位定位的最优观测点分析给出的正是这一权衡关系。

\subsection{几何稀释：Foy 口径}

Foy~\cite{foy1976} 从几何稀释的角度指出：两条示向度夹角越小，源位置沿角平分线方向的误差被
$1/\sin\gamma$ 放大；并且在两条极端射线之间\textbf{再多加测线也不能改善}——必须真正扩大张角。
把式~\eqref{eq:gdop} 中的 $\sigma$ 换成有界误差对应的横向位移因子 $\tan\varepsilon$，即得该
口径的最小外接圆半径

\begin{equation}
  \rho_{\mathrm{MEC}}=\frac{\tan\varepsilon\sqrt{R_1^2+R_2^2}}{\sin\gamma},
  \label{eq:foy}
\end{equation}

它与式~\eqref{eq:gdop} 同构，差别只在 $\sigma$ 与 $\tan\varepsilon$ 的标定。本文的一阶解析式
\eqref{eq:closed} 与它们同以 $1/\sin\gamma$ 为主导，§\ref{sec:res} 的图~\ref{fig:criteria}
把四种口径画在同一张图里，交点与趋势完全一致。

\subsection{与本文集员判据的结构性差别}

式~\eqref{eq:gdop}、\eqref{eq:crlb}、\eqref{eq:foy} 与本文式~\eqref{eq:closed} 的差别集中在
根号内的\textbf{交叉项}：

\begin{equation}
  \underbrace{\sqrt{\,R_1^2+R_2^2+2R_1R_2|\cos\gamma|\,}}_{\text{集员（本模型最坏）}}
  \qquad\text{对}\qquad
  \underbrace{\sqrt{\,R_1^2+R_2^2\,}}_{\text{统计（独立误差期望）}}.
  \label{eq:cross}
\end{equation}

前者源于"有界误差\textbf{可以同向叠加}"（两个 $\pm1^\circ$ 的有界误差都朝同一侧偏时，
两条射线的位移相干叠加，这正是最坏情况必须承担的）；后者源于"独立误差\textbf{平方向加}"
（交叉项在期望意义下相消）。两者的标定是同一的（都以 $\varepsilon$ 作为角度尺度，
式~\eqref{eq:gdop} 中取 $\sigma=\varepsilon$），因此可以直接比较数值。

\begin{table}[H]
  \centering
  \caption{本模型最坏情形（$d=1500$ m、$r_2=907.24$ m、$\gamma=40.64^\circ$）下各口径给出的
  定位区域半径。前两行为\textbf{精确构造}与\textbf{本文集员闭式}，后三行为\textbf{文献口径}。}
  \label{tab:criteria}
  \small
  \setlength{\tabcolsep}{5pt}
  \begin{tabular}{lccl}
    \toprule
    \textbf{口径} & \textbf{半径 / m} & \textbf{相对偏差} & \textbf{口径} \\
    \midrule
    精确构造（$\pm1^\circ$ 楔形交，本文判据） & $67.04$ & ---               & 有界误差最坏界 \\
    本文集员闭式（式~\eqref{eq:closed}）      & $60.74$ & $-9.4\%$          & 有界误差一阶近似 \\
    文献 GDOP（式~\eqref{eq:gdop}）          & $46.98$ & $-29.9\%$         & 统计 $1\sigma$ 尺度 \\
    文献 CRLB 主半轴（式~\eqref{eq:crlb}）   & $44.77$ & $-33.2\%$         & 统计 $1\sigma$ 尺度 \\
    文献几何稀释式（式~\eqref{eq:foy}）      & $46.98$ & $-29.9\%$         & 有界但独立误差口径 \\
    \midrule
    文献 CRLB 椭圆面积                       & $1999.57$ m$^2$ & ---     & 面积口径，不与半径比较 \\
    \bottomrule
  \end{tabular}
\end{table}

表~\ref{tab:criteria} 说明：\textbf{文献判据在本模型上系统性偏低约三成}。原因不是文献式写错了，
而是口径不同：它们给出的是"独立随机误差下 $1\sigma$ 的均方尺度"，而本文必须给出"有界误差下
最坏情形的硬界"。若把文献式乘以覆盖率因子（例如 $3\sigma$ 对应 $99.7\%$），数值上当然可以
超过本文的界，但那需要引入误差分布假设与置信水平；本模型不假设分布，因此直接取硬界。
清除判据是 $20$ m 的硬阈值（要求"保证成功"而不是"平均够准"），用低估三成的量去保证是不可
接受的，这就是本文\textbf{以集员最坏界作为选点判据}的理由。

抽样的偏差分布进一步说明了适用范围。表~\ref{tab:buckets} 给出 $300$ 例随机几何（固定种子）
按交会角与距离比分桶的相对偏差：本文闭式在 $\gamma\ge15^\circ$ 且 $r_2/d\lesssim5$ 时中位
偏差在 $0.2\%$ 以内（可用于快速估算），但在极长条构型（$r_2/d>5$）会严重偏离；而文献
GDOP 类口径在各桶上稳定偏低 $5\%\sim25\%$。这恰好界定了两者的用途：本文闭式是"同一保守
口径下的解析近似"，文献式是"不同口径下的趋势工具"。

\begin{table}[H]
  \centering
  \caption{两种闭式的偏差分桶（$300$ 例抽样，$n$ 为有效样本数；相对精确最坏半径的中位偏差）}
  \label{tab:buckets}
  \begin{tabular}{lrrr}
    \toprule
    \textbf{交会角 $\gamma$} & $n$ & \textbf{本文集员闭式（最大）} & \textbf{文献 GDOP 式中位} \\
    \midrule
    $[15^\circ,30^\circ)$   & $31$  & $-0.79\%$（$4.4\%$）   & $-25.1\%$ \\
    $[30^\circ,60^\circ)$   & $60$  & $-0.19\%$（$0.4\%$）   & $-19.5\%$ \\
    $[60^\circ,90^\circ)$   & $38$  & $-0.07\%$（$0.1\%$）   & $-7.4\%$ \\
    $[90^\circ,120^\circ)$  & $17$  & $+0.02\%$（$1.3\%$）   & $-1.6\%$ \\
    $[120^\circ,155^\circ)$ & $15$  & $+0.02\%$（$36.0\%$）  & $-10.0\%$ \\
    \bottomrule
  \end{tabular}
\end{table}

\subsection{判据一致性：为什么可以用它做快筛}

虽然文献口径\textbf{定量}偏低，但它的\textbf{量级排序}是否正确？本文在可行域内等距抽取
$1276$ 个可行点，比较两者对同一候选点的排序（Spearman 秩相关）：

\begin{equation}
  \rho_{\mathrm{spearman}}\bigl(\mathrm{GDOP},\,J\bigr)=0.998.
  \label{eq:spearman}
\end{equation}

排序几乎一致——这意味着可以用解析 GDOP 场在细网格上\textbf{快速筛出少量候选}（解析场比精确
构造快两个数量级），再对候选做精确复核与细化。图~\ref{fig:criteria}(b) 给出一致性散点图。
"标的排序一致、绝对标定偏低"正是"可作预筛、不可作保险界"的典型情形，本文据此设计了
§\ref{sec:algo} 的全局认证算法；认证结果与粗搜解仅差 $8.5$ m，两条独立路线互相印证。

最后需要说明文献对本模型的\textbf{定性}支撑，这些结论与本文的几何直觉完全一致：
Foy~\cite{foy1976} 的稀释效应解释了式~\eqref{eq:closed} 的分母 $\sin\gamma$，也解释了为什么
"再多的同侧测线无益"；Chan 与 Ho~\cite{chan1994} 指出在坏 GDOP 几何下最小二乘类定位方法会显著退化
甚至发散，这与本文"必须避开 $\gamma\to0$ 的退化构型"的结论同向；Chen 等~\cite{chen2009} 指出观测点
优化是"靠近目标"与"改善方位分布"的权衡，与本文"两个约束同时顶到边界"的结构完全对应；
任叶童~\cite{ren2016} 的误差椭圆式则提供了本文用于对照的闭式来源。
